
\chapter{Ablations and Failure Analysis}
\label{ch:ablations}

This chapter isolates the contribution of each DC3S component through
systematic ablation and identifies the conditions under which individual
components weaken or fail.  The analysis confirms that safety is preserved
even when individual sub-modules degrade, because the shield's conservatism
compensates for upstream estimation errors.

\section{Reliability Scoring Ablation}

The Observation Quality Engine (OQE) produces a scalar reliability score
$w_t \in [0,1]$.  When the OQE is ablated (replaced by $w_t \equiv 1$),
the shield loses its ability to adapt uncertainty width to telemetry
quality.

\begin{table}[ht]
\centering
\caption{Effect of OQE ablation on DC3S safety metrics (\texttt{drift\_combo}, DE).}
\label{tab:oqe-ablation}
\begin{tabular}{lrrr}
\toprule
\textbf{Configuration} & \textbf{TSVR (\%)} & \textbf{Sev.\ P95 (MWh)} & \textbf{IR (\%)} \\
\midrule
Full DC3S             & 0.0 &  0.0 & 2.8 \\
OQE ablated ($w=1$)   & 0.0 &  0.0 & 0.4 \\
OQE ablated ($w=0.5$) & 0.0 &  0.0 & 8.3 \\
\bottomrule
\end{tabular}
\end{table}

With $w_t \equiv 1$ the shield rarely intervenes (0.4\%) because it
trusts all observations, yet safety is preserved on this scenario because
the conformal intervals alone are wide enough.  However, under more
aggressive faults the fixed-$w$ controller would eventually miss violations.

\section{Inflation Ablation}

The Reliability-aware Uncertainty Inflation (RUI) module scales conformal
intervals by $1/(w_t + \epsilon)$.  When inflation is disabled (intervals
are used at their nominal CQR width), the shield under-covers the true
state under high dropout:

\begin{table}[ht]
\centering
\caption{RUI ablation under varying dropout rates (DE).}
\label{tab:rui-ablation}
\begin{tabular}{lrrrr}
\toprule
\textbf{Dropout} & \multicolumn{2}{c}{\textbf{Full DC3S}} &
  \multicolumn{2}{c}{\textbf{RUI ablated}} \\
\cmidrule(lr){2-3} \cmidrule(lr){4-5}
& TSVR & IR & TSVR & IR \\
\midrule
0.05 & 0.0 & 0.8 & 0.0 & 0.6 \\
0.15 & 0.0 & 1.9 & 0.2 & 0.9 \\
0.25 & 0.0 & 2.6 & 1.1 & 1.2 \\
\bottomrule
\end{tabular}
\end{table}

RUI ablation reveals non-zero TSVR at dropout rates above 0.10,
confirming that reliability-adaptive inflation is necessary for the
zero-violation guarantee.

\section{Projection Ablation}

The Safe-Action Feasibility (SAF) projection maps infeasible actions onto
the tightened constraint set.  When projection is replaced by simple
clipping (independently clamp each variable to its bound), the solver
can violate coupling constraints between charge and discharge power:

\begin{itemize}
  \item Full projection: TSVR~$= 0$\%, IR~$= 2.8$\%.
  \item Clipping only: TSVR~$= 0.3$\%, IR~$= 2.1$\%.
\end{itemize}

The clipping ablation shows that joint constraint projection (the LP step)
is required to maintain zero violations when charge/discharge variables
are coupled by ramp and SOC dynamics.

\section{FTIT Ablation}

The Forecast-Triggered Interval Tightening (FTIT) module narrows the
conformal interval when recent forecast errors are small, preventing
unnecessary conservatism in calm periods.  Disabling FTIT does not
affect safety but increases the intervention rate:

\begin{itemize}
  \item Full DC3S: IR~$= 2.8$\%.
  \item FTIT disabled: IR~$= 4.6$\%.
\end{itemize}

FTIT is therefore a performance refinement, not a safety-critical component.

\section{Fault-Family Sensitivity}

Table~\ref{tab:fault-sensitivity} decomposes TSVR by fault family for the
deterministic LP baseline, revealing which telemetry failures cause the
most damage.

\begin{table}[ht]
\centering
\caption{Deterministic LP TSVR by individual fault family (DE, dropout~0.20).}
\label{tab:fault-sensitivity}
\begin{tabular}{lr}
\toprule
\textbf{Fault family} & \textbf{TSVR (\%)} \\
\midrule
Packet dropout only       & 2.1 \\
Delay jitter only         & 0.8 \\
Covariate drift only      & 3.4 \\
Stale sensor only          & 1.6 \\
Combined (\texttt{drift\_combo}) & 3.9 \\
\bottomrule
\end{tabular}
\end{table}

Covariate drift is the most damaging single fault, consistent with the
observation that drift induces systematic bias rather than zero-mean noise.

\section{When DC3S Intervenes}

Interventions cluster in three regimes:
\begin{enumerate}
  \item \textbf{Peak-load transitions} where forecast error spikes.
  \item \textbf{Consecutive dropout bursts} (3+ consecutive dropped packets).
  \item \textbf{Drift-ramp events} where covariate shift exceeds the CQR
    calibration window.
\end{enumerate}
Outside these regimes, the shield is transparent: $a_t^{\text{safe}} = a_t^{\text{LP}}$.

\section{When Base Calibration Weakens}

The conformal calibration set is a trailing window of 500 residuals.
Calibration weakens when:
\begin{itemize}
  \item The distribution shifts faster than the window refreshes.
  \item The miscoverage rate temporarily exceeds $\alpha$ during a
    transient regime change.
\end{itemize}
Empirically, the CQR marginal coverage dips to 88\% during the first
30~minutes of an abrupt drift onset (nominal target: 90\%).

\section{Why Safety Is Preserved Anyway}

Even when calibration weakens, safety is preserved because:
\begin{enumerate}
  \item The RUI module inflates the interval proportionally to the
    reliability drop, compensating for the coverage gap.
  \item The SAF projection adds a second line of defence by mapping
    any remaining infeasible action onto the constraint set.
  \item The reliability score $w_t$ drops faster than coverage degrades,
    triggering pre-emptive tightening before violations can occur.
\end{enumerate}

This layered defence is the operational manifestation of the theoretical
safety-preservation theorem (Theorem~\ref{thm:safety-preservation}).

\section{Rank Reversal Under Degradation}

A key finding from the CPSBench sweep is that controller rankings are not
stable across fault conditions.  Table~\ref{tab:rank_reversal} shows how
each controller's rank changes from the nominal (clean telemetry) to the
degraded (30\% dropout, simultaneous drift) scenario.  CVaR Interval
drops by 5 positions; DC3S-FTIT maintains or improves its rank because
its reliability-adaptive inflation correctly widens the uncertainty set
rather than assuming clean inputs.

\begin{table}[htbp]
\centering
\caption{Rank Reversal Under Telemetry Degradation: nominal vs.\ degraded scenario rankings.  Positive rank shift = improvement under stress.}
\label{tab:rank_reversal}
\begin{tabular}{lrrrr}
\toprule
Controller & Nominal Rank & Degraded Rank & Rank Shift & Degraded Score \\
\midrule
CVaR Interval & 1 & 7 & -6 & 0.2 \\
Robust Fixed & 1 & 7 & -6 & 0.2 \\
Deterministic LP & 3 & 6 & -3 & 0.1 \\
Scenario MPC & 4 & 1 & +3 & 0.0 \\
Scenario Robust & 4 & 1 & +3 & 0.0 \\
ACI Conformal & 6 & 1 & +5 & 0.0 \\
DC3S (Wrapped) & 7 & 1 & +6 & 0.0 \\
DC3S-FTIT & 8 & 1 & +7 & 0.0 \\
\bottomrule
\end{tabular}
\end{table}

\section{Chapter Summary}

The ablation study confirms that each DC3S component contributes to
either safety or efficiency, but the zero-violation property depends
critically on the combination of reliability-aware inflation (RUI) and
joint constraint projection (SAF).  Neither component alone is sufficient
under high-dropout combined faults.  The shield's layered architecture
ensures that temporary calibration weaknesses are absorbed before they
propagate to true-state violations.
